% =========================================================================
% 摘要页（front matter）
% =========================================================================

\pagenumbering{arabic}
\begin{mcmabstract}
  本文研究微网与外部电网电力调控中的储能调度与购电策略。四问共用同一储能物理模型（容量 12000 kWh、允许储电量 $[1200,10800]$ kWh、充放电功率上限 5000 kW、双向效率 90\%），差别在于决策时可用的信息与结算规则：问题一为给定输入下的确定性日调度；问题二引入未来负载与光伏的未知性，缺口按 5 倍电价惩罚；问题三在问题二之上增加 6:00、12:00、18:00 三个预报时刻，并允许付费重定合同；问题四把电价本身变为逐区间波动的实时价。全文以``统一物理层＋各问新增信息与决策权限＋受控对照''为主线，目标为最小化购电费用。

  针对\textbf{问题一}，将一天离散为 144 个 10 分钟区间，建立以全天购电费用最小为目标的线性规划，并证明充放电互斥约束的线性松弛在最优值意义下精确，从而无需引入整数变量。以 HiGHS 求得的策略为：全天购电费 \textbf{35126.95 元}、购电量 59482.70 kWh，较无储能基准（48052.05 元）\textbf{节省 26.90\%}。

  针对\textbf{问题二}，构造``预测—计划—执行''三层滚动框架：预测层以近期净负载水平、星期偏差校正与分位裕量三项叠加给出当日预测；计划层求解带日循环边界的日前线性规划；执行层按净缺口符号即时尽限充放电，不足部分按 5 倍电价紧急购电。预测器形式、窗口长度与裕量水平均按逐月 walk-forward 协议选取，严格排除数据泄漏。评价期（2025.2.1--12.31，共 334 天）总费用 \textbf{1395.70 万元}，其中紧急费用占 5.5\%。

  针对\textbf{问题三}，在问题二框架上增设\textbf{调整层}：0:00 用附件 3 预报经分档收缩回归修正净负载预测，6:00、12:00、18:00 在实测储电量下对剩余时段重解带双向违约价（下调 $0.5p$、上调超出部分 $1.5p$）的线性规划，并证明该含分段的结算规则可由\emph{单个}线性规划精确表示。以``是否使用 0:00 预报通道 $\times$ 是否允许日内调整''构造 $2\times2$ 对照，主方案总费用 \textbf{1345.71 万元}，较问题二降费 \textbf{49.99 万元}，两项改进存在 9.62 万元的重叠。对题面所问``是否需要引入其他时刻的预报''，在本文所比较的全融合与逐点门控两种策略下，日内新增预报均未进一步降低费用（分别增费 9.59 与 6.08 万元）——双向违约价削弱了预测精度向费用节省的转化渠道。

  针对\textbf{问题四}，在附件 4 的逐区间实时价下重算问题二与问题三两条链路。价格的结构分析表明，实时电价接近``分时形状固定、日水平浮动再叠加小残差''的形式；据此可得\textbf{价格缩放不变性}，即价格的整体水平不产生决策价值，决策只依赖其归一化形状。模型层面引入两项修改：按价格形状加权的\textbf{价格加权风险裕量}，以及把优化终点延至次日 24:00 的 \textbf{48 小时滚动时域}（次日仅用持久性预测，信息严格因果）。两条分支的主方案费用分别为 \textbf{1438.23 万元}与 \textbf{1404.24 万元}；48 小时时域使两分支分别降费 \textbf{11.08 万元}与 \textbf{6.70 万元}，并在共同终端条件下保持。三个不含未来信息的在线选择器均未超过同年探索性选定的固定参数。
\end{mcmabstract}

\clearpage
